%! Inverse Functions — Unit 2, Lesson 2.8 — Warm-Up (Answer Key)
\documentclass[10pt]{article}
\usepackage{precalculus-article}
\usepackage{precalculus-key}   % pulls in -boxes; do NOT also load -boxes
\newcommand{\TallMath}[1]{$\displaystyle #1\rule[-1.4em]{0pt}{3.2em}$}

\begin{document}

\pageheader{Unit 2, Lesson 2.8}{Warm-Up --- Answer Key}
\namedateperiod

\vspace{0.14in}

\begin{enumerate}[itemsep=18pt]

\item Let $f(x) = 2x + 3$. Find the value of $x$ when $f(x) = 11$. Show your work.

      \vspace{0.06in}
      $2x + 3 = 11$

      \vspace{0.04in}
      $x = $ \ans{4}

\item Suppose $f(g(x)) = x$ for all $x$ and $g(f(x)) = x$ for all $x$.
      What is the mathematical relationship between $f$ and $g$?

      \vspace{0.06in}
      \ansline{$f$ and $g$ are inverse functions of each other; each undoes the other.}
      \ansline{We write $g = f^{-1}$ (or equivalently $f = g^{-1}$).}

\item A function $f$ maps: $1 \to 5$, $2 \to 8$, $3 \to 11$.

      If we reverse all input-output pairs, what does the new function map $8$ to?

      \vspace{0.06in}
      The new function maps $8$ to: \ans{2}

\end{enumerate}

\end{document}
